\paragraph{What the diagnostics reveal.}
The central lesson from our experiments is that long-horizon route generation fails for a structural reason: corridor-level multi-modality is \emph{topological} and cannot be reliably encoded as a single mode in a continuous coordinate space.
This manifests as destructive averaging in deterministic regression, and as corridor-level interference in stochastic continuous generators.

\paragraph{Why hierarchy is necessary but not sufficient.}
Separating decision from execution resolves the representation mismatch by letting the model commit to discrete corridor alternatives.
However, our results also reveal a subtle failure mode: if execution is too expressive, it can overwrite the commitment and reintroduce corridor collapse.
This suggests that hierarchical route generators should be designed with an explicit ``commitment preservation'' interface, rather than treating the executor as a generic conditional generator.

\paragraph{Feasibility and map information.}
We view map-derived information as a useful source of inductive bias but not a definition of support.
Hard masking can enforce on-road outputs but risks truncating the learned distribution and entangling performance with proxy quality.
Soft priors (e.g., road proximity fields) offer a spectrum of feasibility guidance while keeping the generator distributional.

\paragraph{Limitations and scope.}
This draft emphasizes a diagnostic evidence chain on fixed multi-modal OD cases; scaling to broader city-wide benchmarks and reporting full realism/feasibility tradeoffs are ongoing.
Our core findings do not require external semantic data; nevertheless, the decision stage offers a natural interface for optional spatial covariates (e.g., POI/land-use or census attributes).
We find that coarse OD-point semantics can improve corridor-level diversity and mixture matching, while naive corridor-agnostic ``strong'' semantic summaries may hurt, indicating that exploiting rich semantics may require corridor-aware representations (e.g., candidate corridor evaluation) rather than a single global profile.
We do not attempt full population synthesis or attribute generation in this paper.
